\exercice*

À chaque question de cet exercice, on utilise le fait que la courbe d'une fonction $(( f ))$ est, par définition, l'ensemble des points de coordonnées $\left(x; (( f ))(x)\right)$, $x$ appartenant à l'ensemble de définition de $(( f ))$.

On rappelle que la fonction $((f))$ est définie sur $\mathbb{R}$ par :
(* if nature ==  "degre1" *)
$(( f ))(x) = (( a|facteur("*x") ))(( b|facteur("so") ))$.
(* elif nature ==  "degre2" *)
$(( f ))(x) = (( a|facteur("*X") ))(( b|facteur("*sox") ))(( c|facteur("so") ))$.
(* endif *)

\begin{enumerate}
  (* if sens ==  "antecedent-vers-image" *)
    \item Tous les points de l'axe des ordonnées ont pour abscisse 0. Les coordonnées du point recherché sont donc $\left(0; (( f ))(0)\right)$. Puisque 
        (* if nature ==  "degre1" *)
        $(( f ))(0)=(( a ))\times0 (( b|facteur("so") ))=(( b ))$,
        (* elif nature ==  "degre2" *)
        $(( f ))(0)=(( a ))\times0^2 (( b|facteur("so") ))\times0 (( c|facteur("so") ))=(( c ))$,
        (* endif *)
        alors les coordonnées du point d'intersection entre la courbe de $(( f ))$ et l'axe des ordonnées sont $\left( 0; (( C.y ))\right)$.
  (* else *)
  \item Tous les points de l'axe des abscisses ont pour ordonnée 0. Les coordonnées du point recherché sont donc de la forme $\left( x; 0 \right)$, où $x$ est à déterminer. Mais puisque les coordonnées des points de la courbe de $(( f ))$ sont $\left( x;(( f ))(x) \right)$, alors nous recherchons $x$ tel que $(( f ))(x)=0$ :
    \begin{align*}
        (( f ))(x) &= 0 \\
        (( a|facteur("*x") )) (( b|facteur("so") )) &= 0\\
        (( a|facteur("*x") )) &= (( -b )) \\
        x &= \frac{(( -b ))}{((a))} \\
        (* if a|pgcd(b) != 1 *)
        x &= (( Cxfinal ))
        (* endif *)
    \end{align*}
    Donc il y a un unique point d'intersection entre la courbe de $(( f ))$ et l'axe des abcsisse, qui a pour coordonnées $\left( ((Cxfinal));0 \right)$.
  (* endif *)
\item Ce point $A( ((A.x)) ; ((A.y)) )$ est sur la courbe de $(( f ))$ si et seulement si $(( f ))( ((A.x)) )=((A.y))$. Vérifions cela :
    \[
        (* if nature ==  "degre1" *)
        (( f ))( ((A.x)) ) = (( a ))\times(( A.x|facteur("p") )) (( b|facteur("so") ))=(( a * A.x )) ((b|facteur("so"))) = (( a*A.x+b ))
        (* else *)
        (( f ))( ((A.x)) )
        = (( a ))\times(( A.x|facteur("p") ))^2 (( b|facteur("so") ))\times (( A.x|facteur("p") )) (( c|facteur("so") ))
        = (( a ))\times(( A.x**2 )) (( (b*A.x)|facteur("so") )) (( c|facteur("so") ))
        = (( a*A.x**2 )) (( (b*A.x)|facteur("so") )) (( c|facteur("so") ))
        = (( a*A.x**2 + b*A.x + c ))
        (* endif *)
    \]
    (* if 
        ( nature == "degre1" and a*A.x+b == A.y )
        or
        ( nature == "degre2" and a*A.x**2 + b*A.x + c == A.y )
    *)
    Donc $(( f ))( ((A.x)) )=((A.y))$, et le point $A( ((A.x)); ((A.y)) )$ est sur la courbe de $(( f ))$.
    (* else *)
    Donc $(( f ))( ((A.x)) )\neq((A.y))$, et le point $A( ((A.x)); ((A.y)) )$ n'est pas sur la courbe de $(( f ))$.
    (* endif *)
    (* if sens ==  "antecedent-vers-image" *)
        \item Ce point $B( ((B.x)) ;y)$ est sur la courbe de $(( f ))$, donc :
            (* if nature ==  "degre1" *)
            $y=(( f ))( ((B.x)) ) = (( a ))\times(( B.x|facteur("p") )) (( b|facteur("so") ))=(( a * B.x )) ((b|facteur("so"))) = (( B.y ))$.
            (* elif nature ==  "degre2" *)
            $y=(( f ))( ((B.x)) )
            = (( a ))\times(( B.x|facteur("p") ))^2 (( b|facteur("so") ))\times (( B.x|facteur("sop") )) (( c|facteur("so") ))
            = (( a ))\times(( B.x**2 )) (( b|facteur("so") ))\times (( B.x|facteur("sop") )) (( c|facteur("so") ))
            = (( B.y ))$.
            (* endif *)
            La valeur recherchée est $y=(( B.y ))$.
        (* else *)
    \item Le point $B(x;((B.y)) )$ est sur la courbe de $(( f ))$ si et seulement si $(( f ))(x)=((B.y))$. Résolvons cette équation pour trouver les valeurs possibles de $x$.
    \begin{align*}
        (( f ))(x)&= ((B.y))\\
        (( a|facteur("*x") )) (( b|facteur("so") )) &= ((B.y)) \\
        (( a|facteur("*x") )) &= (( B.y ))(( (-b)|facteur("so") )) \\
        (( a|facteur("*x") )) &= (( B.y-b )) \\
        x &= \frac{(( B.y-b))}{(( a ))}\\
        x &= (( B.x ))
    \end{align*}
    Il y a donc une seule solution possible : $x=(( B.x ))$.
    (* endif *)
\end{enumerate}
